\documentclass[11pt]{amsart}
\usepackage[T1]{fontenc}
\usepackage{amsmath,amssymb,amsthm}
\usepackage[colorlinks=true,linkcolor=blue,citecolor=blue,urlcolor=blue]{hyperref}

\newtheorem*{theorem*}{Literature-implied answer}

\begin{document}

\title{A One-Test-Space Characterization for Arbitrary Finite-Dimensional Sequences}
\author{Run fa\_banach\_001}
\date{2026-06-19}
\maketitle

\noindent
\textbf{Status.} \texttt{literature\_implied\_answer (full arbitrary-sequence
one-test-space, bounded-degree graph)}. This is not presented as a new theorem
of the run. It records that a later finite-test-space-to-one-test-space
proposition, combined with the finite graph test-space theorem in the source
paper, gives the arbitrary-sequence case that was not known in the source's
Section 3.

\section*{Source Question}

Ostrovskii's \emph{Test-space characterizations of some classes of Banach
spaces}, arXiv:1112.3086, defines test-spaces for a class \(\mathcal P\) of
Banach spaces by the equivalence
\[
  X\notin \mathcal P
  \quad\Longleftrightarrow\quad
  \text{the test-spaces admit uniformly bilipschitz embeddings into } X.
\]
Theorem \(\mathrm{T:Ost11+}\) of that paper states that for every sequence
\((X_m)_{m=1}^{\infty}\) of finite-dimensional Banach spaces there is a
sequence \((H_n)_{n=1}^{\infty}\) of finite connected unweighted graphs, each
of maximum degree \(3\), such that for every Banach space \(Y\),
\[
  (X_m) \text{ admit uniformly isomorphic embeddings into } Y
  \quad\Longleftrightarrow\quad
  (H_n) \text{ admit uniformly bilipschitz embeddings into } Y.
\]

In Section 3, after recalling one-test-space characterizations for
superreflexivity, the source paper says that it seeks analogous
one-test-space characterizations for classes defined by excluded
finite-dimensional subspaces, but that at that moment this was known only for
increasing sequences, not for arbitrary sequences.

\section*{Supporting Literature}

Ostrovskii's later \emph{Metric characterizations of some classes of Banach
spaces}, arXiv:1411.3366, records the following elementary proposition in the
discussion of one test-spaces. If \((S_n)\) is a sequence of finite test-spaces
for a class \(\mathcal P\) of Banach spaces containing all finite-dimensional
Banach spaces, then there is a single metric space \(S\) which is a test-space
for \(\mathcal P\). Moreover, if the \(S_n\) are unweighted graphs, trees, or
graphs with uniformly bounded degrees, then \(S\) may be required to have the
same corresponding graph property.

The proof sketch constructs \(S\) by choosing vertices \(O_n\in S_n\) and
joining \(O_n\) to \(O_{n+1}\) by long paths. The nontrivial direction embeds
the \(S_n\)'s in parallel hyperplanes of the target Banach space, separated by
the path lengths.

\begin{theorem*}
Let \((X_m)_{m=1}^{\infty}\) be an arbitrary sequence of finite-dimensional
Banach spaces with \(\sup_m \dim X_m=\infty\). Then there is a single metric
space \(S\) such that, for every Banach space \(Y\),
\[
  S \text{ admits a bilipschitz embedding into } Y
  \quad\Longleftrightarrow\quad
  (X_m) \text{ admit uniformly isomorphic embeddings into } Y.
\]
Furthermore, \(S\) can be chosen to be an unweighted graph with uniformly
bounded degree.
\end{theorem*}

\begin{proof}
Let
\[
  \mathcal P=\{Y:\ (X_m) \text{ do not admit uniformly isomorphic embeddings
  into } Y\}.
\]
Since \(\sup_m\dim X_m=\infty\), every finite-dimensional Banach space belongs
to \(\mathcal P\): no finite-dimensional target can linearly contain all
members of an unbounded-dimensional sequence.

Apply Theorem \(\mathrm{T:Ost11+}\) from arXiv:1112.3086 to obtain finite
connected unweighted graphs \((H_n)\), with maximum degree \(3\), such that
\[
  Y\notin \mathcal P
  \quad\Longleftrightarrow\quad
  (H_n) \text{ admit uniformly bilipschitz embeddings into } Y.
\]
Thus \((H_n)\) is a sequence of finite test-spaces for \(\mathcal P\).

The proposition from arXiv:1411.3366 applies because \(\mathcal P\) contains
all finite-dimensional Banach spaces. It yields a single metric space \(S\)
which is a test-space for \(\mathcal P\). Therefore, for every Banach space
\(Y\),
\[
  S \text{ embeds bilipschitzly into } Y
  \quad\Longleftrightarrow\quad
  Y\notin\mathcal P
  \quad\Longleftrightarrow\quad
  (X_m) \text{ admit uniformly isomorphic embeddings into } Y.
\]
Since the \(H_n\)'s are unweighted graphs of uniformly bounded degree, the
same proposition permits \(S\) to be chosen as an unweighted graph with
uniformly bounded degree.
\end{proof}

\section*{Scope}

This answers the source paper's arbitrary-sequence one-test-space gap under
the same nontrivial unbounded-dimension convention used in the source paper's
Remark \(\mathrm{R:InfDim}\). The supporting proposition gives a uniformly
bounded-degree graph. This note does not separately establish a maximum-degree
\(3\) refinement for the arbitrary-sequence one-graph construction.

\section*{Search and Provenance Notes}

The run indexes were searched for arXiv:1112.3086, the title of the source
paper, one-test-space terminology, arbitrary finite-dimensional sequences,
finite-dimensional Banach test-spaces, locally finite and coarse-disjoint
union language, and Ostrovskii. No prior packet for this exact identification
was found. Web queries on June 19, 2026 for exact combinations of
\texttt{1112.3086}, \texttt{one-test-space}, \texttt{arbitrary sequence},
\texttt{finite-dimensional subspaces}, and \texttt{Ostrovskii} found the
source paper but no explicit separate answer. The decisive supporting
proposition was located in the local parsed source of arXiv:1411.3366.

\end{document}
